% TOP_PROBLEM: {"id": 8700053, "title": "Conjecture 4.11 — Wirsing and Schmidt", "category_id": 3, "category": {"id": 3, "name": "graph_theory", "display_name": "Graph Theory", "description": "Problems involving graphs, networks, and their properties.", "slug": "graph-theory", "order_index": 3, "created_at": "2026-07-31T15:26:25.671Z"}, "status": "partially_solved", "problem_number": "AMR-086-0053", "set_id": 13}
% FIRST_POSTED: 2026-10-01
% ENTRY_KIND: statement_only
% UnsolvedMath ID 8700053; source catalogue code AMR-086-0053
% !TeX program = lualatex
% Definition and sources only; this file is not an LLM solution attempt.
\documentclass[11pt,a4paper]{article}
\usepackage[margin=25mm]{geometry}
\usepackage{fontspec}
\setmainfont{lmroman10-regular.otf}[BoldFont=lmroman10-bold.otf,
 ItalicFont=lmroman10-italic.otf,BoldItalicFont=lmroman10-bolditalic.otf]
\usepackage{amsmath,amssymb,mathtools}
\usepackage{unicode-math}
\setmathfont{latinmodern-math.otf}
\AtBeginDocument{\let\setminus\smallsetminus}
\usepackage{xurl,hyperref}
\hypersetup{colorlinks=true,urlcolor=blue}
\setcounter{secnumdepth}{0}
\setlength{\parindent}{0pt}
\setlength{\parskip}{5pt}
\setlength{\emergencystretch}{3em}
\begin{document}
\section*{Conjecture 4.11 — Wirsing and Schmidt}\label{8700053}
{\small UnsolvedMath ID 8700053 · AMR-086-0053 · Graph Theory · AMR Open Problem Lists}
\subsection{Definitions and mathematical statement}
For a real algebraic number $\gamma$, let its primitive irreducible polynomial over $\mathbb Z$ have positive leading coefficient. Its naive height $H(\gamma)$ is the largest absolute value of the coefficients of this polynomial, and its degree is the polynomial's degree. This height differs from the absolute logarithmic Weil height.

Fix an integer $n\ge1$ and a real number $\theta$ which is either transcendental or algebraic of degree strictly greater than $n$. Does there exist a real constant $c=c(n,\theta)>0$ for which infinitely many distinct real algebraic numbers $\gamma$ of degree at most $n$ satisfy
\[
0<|\theta-\gamma|<c\,H(\gamma)^{-n-1}?
\]
The constant is chosen once for the given pair $(n,\theta)$ and must work for the entire infinite collection of approximants. No effective procedure for finding it or the approximants is required.

The degree restriction is an upper bound, so rational approximants and algebraic numbers of any smaller degree are permitted. The assumptions on $\theta$ prevent the trivial repeated approximation $\gamma=\theta$; the strict positive lower inequality also records this exclusion explicitly.

Infinitely many distinct approximants necessarily have unbounded heights, because only finitely many primitive integer polynomials have bounded degree and bounded coefficient height. The target exponent is exactly $n+1$, with a fixed multiplicative constant. Bounds with a smaller exponent, or constants allowed to grow from one approximant to the next, do not establish this formulation.
\subsection{Short English statement}
Can every eligible real number be approximated infinitely often by algebraic numbers of degree at most n at the precise height exponent n+1?
\subsection{Sources}
\begin{itemize}
\item[S1] ULAM.AI and the UnsolvedMath Contributors, supplied problems.json, record 8700053 (AMR-086-0053), AMR Open Problem Lists. Catalogue identity and source transcription; CC BY 4.0.
\url{https://huggingface.co/datasets/ulamai/UnsolvedMath}.
\item[S2] Waldschmidt - Open Diophantine Problems (2004), Conjecture 4.11. Original source indexed by AMR-086-0053.
\url{https://arxiv.org/abs/math/0312440}.
\end{itemize}
\end{document}
